\begin{abstract}
Disputes about linguistic categories often combine different questions about
the same expressions. English interjections provide a worked case.
Grammatical interjections, expressive signs, depictive signs, and particular
interactional resources preserve different relations, including differences
within a single semantic inquiry. I distinguish constructed classes, useful
for measurement and description, from empirical categories whose co-occurring
properties support induction for particular purposes. Each proposal specifies
its bearer, grouping basis, intended contribution, scope, and alternatives.
A comparison of expressive and depictive uses shows what an undifferentiated
interjection judgment obscures; a phonological case shows why evidence for
one grouping cannot settle a grammatical assignment. Historical explanations
and causal assessments require further evidence, with Hill's considerations
applying to specified causal claims. The relation among fields, inquiries,
and categories is many-to-many: neither a discipline nor a question determines
one privileged grouping. This account locates disagreement more precisely
and makes its evidential consequences explicit, without treating descriptive
usefulness as proof of categoryhood or every failed projection as grounds
for abandoning all related distinctions.
\end{abstract}

\noindent\textbf{Keywords:} inquiry-relative categories; category pluralism;
projectibility; interjections; sign-types; path dependence

\aidisclosure{Claude Opus 4.8, Claude Opus 5, Claude Opus 5.5, GPT-5.6 Sol,
GPT-6 Astra Pro, and GPT-6 (Codex)}

\section{Introduction}\label{sec:intro}

Standalone \mention{Damn!}, complement-taking \mention{Damn these mosquitoes!},
and prenominal \mention{damn} in \mention{the damn car} present different
combinations of grammatical and semantic properties. The modifier participates
in a noun phrase while retaining an expressive contribution. The standalone
form can constitute a stance display. The complement-taking expression has
internal structure even when the whole expression stands apart from a
surrounding clause. A single judgment about whether \mention{damn} is an
interjection leaves these differences unresolved.

A comparison with depictive \mention{Bang!} exposes a further distinction.
The selected \mention{damn} use conveys an evaluative stance; the selected
\mention{Bang!} depicts a sound. Both concern meaning, but they motivate
different groupings and different questions about interpretation. Grouping
one under semantics and the other under semiotics would do nothing to
resolve those questions. Nor would calling both interjections establish
that they share the relevant semantic properties.

This paper asks what warrants recognizing a category for an inquiry.
An inquiry identifies a question and a standard of success; evidence
constrains which grouping answers it. A field can pursue several inquiries,
an inquiry can employ several categories, and a category can serve questions
pursued across fields. English interjections make this many-to-many relation
visible because grammatical, semantic, interactional, phonological, and
historical descriptions repeatedly organize overlapping material.

A terminological distinction is essential. Stipulating the membership
conditions of a constructed \term{class} is unproblematic. I reserve
\term{category} for empirical groupings whose co-occurring properties support
induction for particular purposes, following the usage in
\textit{Words That Won't Hold Still} \citep[ch.~2]{reynolds2025hpcbook}.
\term{Classification} assigns objects under a chosen scheme;
\term{categorization}, in this usage, identifies them with an empirical
category. These are local conventions, not terminology imposed on the
sources discussed below.

A class can make measurement comparable or description informative without
yet establishing such a category. The distinctions among \mention{damn}
uses have descriptive value even before a new prediction is tested.
\term{Projectibility}, the warrant for expecting more than was used to
identify a case, is a further achievement \citep{Goodman1955}.
Neither descriptive usefulness nor a successful projection identifies its
own causal explanation. The paper separates these achievements so that
evidence for one isn't made to stand in for another.

Pluralism about useful and objectively grounded groupings has substantial
precedents. \textcite[5--7, 17--18]{dupre1993disorder} relates the choice of
classification to the goals of inquiry while rejecting one uniquely
privileged division of things. \textcite[63--64, 69--72]{khalidi2013}
defends crosscutting kinds answering different epistemic interests.
Their wider metaphysical commitments needn't all be adopted to recognize
that distinct questions can require overlapping groupings.

Within linguistics, \textcite[663, 666, 680]{haspelmath2010} distinguishes
language-particular descriptive categories from concepts designed for
cross-language comparison. The division of labour here is different:
several inquiries within English can organize the same expressions
differently. The contribution is to connect that plurality with explicit
bearers, comparisons among alternative groupings, and separate standards
for description, projection, representation, and causal explanation.

The worked cases are named for the relations they investigate: grammatical
interjections, expressive signs, depictive signs, continuers, repair
initiators, and closing formats. They needn't all operate at the same level
of generality or retain an interjection label. In these comparisons, the
common bearer is a conventionalized sign-type, an identifiable expression
in a conventional use. Classifications of realizations or historical
changes require explicit relations to these sign-types.

The historical case concerns \term{path dependence}: development can leave
source-specific traces after uses converge. Gehweiler's history of
\mention{gee} supplies a documented explanation of one outcome; a verbal
route to complement-taking \mention{damn} remains partly reconstructive.
The corpus material has a narrower role, showing how tags and strings
can fail to identify the intended bearers. A linked supplement preserves
the detailed measurement checks and two withdrawn confirmatory analyses.

Section~\ref{sec:worked} develops the named groupings and completes an
expressive--depictive comparison. Section~\ref{sec:assessment} distinguishes
their possible achievements. Sections~\ref{sec:phonology} and~\ref{sec:history}
apply the assessment to phonological and historical claims;
Section~\ref{sec:extensions} considers broader social and processing
classes. The measurement case and discussion then draw out what changes
when the object, grouping, and claim are kept explicit.

\section{Several groupings of conventional signs}\label{sec:worked}

Grammars and studies of meaning and interaction organize the traditional
interjection domain differently. \textcite[90]{jespersen1924} treated
interjection as a manner in which words of other categories are used.
\textit{The Cambridge Grammar of the English Language} (\textit{CGEL})
recognizes a lexical category \citep[22]{HuddlestonPullum2002}. The
\textit{Longman Grammar of Spoken and Written English} places interjections
within the broader class of inserts, while assigning \mention{damn},
\mention{gosh}, and \mention{gee} to expletives
\citep[1082--1083, 1094--1095]{Biber1999}.

The disputes concern a well-studied domain. \textcite{libert2019} questions
the description of interjections as neglected; \textcite[63]{cram2008}
identifies their recurring role in controversies over the syntax--pragmatics
and semantics--pragmatics boundaries. The following comparisons preserve
what those traditions distinguish without treating each tradition as the
source of exactly one category.

\subsection{Grammatical interjections and expressive signs}

The grammatical proposal begins with standalone and supplement distribution.
In \mention{Damn, we're going to be late!}, \mention{damn} is appended to
the clause without being its complement or modifier. \textit{CGEL} treats
interjections as occurring on their own or functioning as supplements
\citep[1350, 1361]{HuddlestonPullum2002}. In those conventional uses,
they're characteristically invariant and resist the ordinary constituent
positions available to nouns, verbs, or adjectives.

Supplement function is shared by several categories. Compare the interjection
with the relative clause and noun phrase in these examples, adapted from
\textit{CGEL} \citep[1356--1357, 1360]{HuddlestonPullum2002}.

\ea\label{ex:supplements}
  \ea \exmention{\exfocus{Ah}, so you were there after all!}
  \ex \exmention{We called in to see Sue's parents, \exfocus{which made us rather late}.}
  \ex \exmention{A university lecturer, \exfocus{Dr Brown}, was arrested for the crime.}
  \z
\z

Their external relation to a host leaves their internal grammatical
differences intact. The same applies to \mention{Damn these mosquitoes!}:
the whole phrase's detachment is compatible with complementation inside it.
\textit{CGEL} tentatively treats the head as an exceptional interjection
taking a noun-phrase complement \citep[1361 n.~70]{HuddlestonPullum2002}.
A verbal analysis remains a substantive alternative. Neither the
complement nor the whole phrase's independence decides between them alone.

Prenominal \mention{damn}, in \mention{the damn car}, occurs inside the
noun phrase while contributing an attitude toward its referent.
\textcite[171, 182]{potts2007} distinguishes that expressive contribution
from descriptive content. Grammatical integration therefore leaves a
semantic relation to standalone \mention{Damn!} available for analysis.
A single degree of interjectionhood would conceal which respect is shared
and which differs.

The semantic grouping here is \term{expressive signs}: conventional uses
whose deployment conveys an evaluative or affective stance, rather than
merely describing that stance. This intended contribution is narrower than
all non-descriptive or use-conditional meaning. A conventional instruction
for managing interpretation or interaction needn't express such a stance.
Potts's diagnostics include perspective dependence, immediacy, and
intensifying repetition \citep[166--167]{potts2007}; whether they travel
together across a proposed extension is an empirical question.

The grouping is broader than grammatical interjections. The examples in
this paper are a selected part of it, chosen from disputes about
interjections; that selection supplies no additional semantic membership
condition. A broader account of expressives can include integrated signs
with no traditional interjection label. Its descriptive benefit here is
to preserve the expressive--descriptive distinction across independently
identified grammatical constructions.

\subsection{Expressive and depictive signs: a completed comparison}\label{sec:expressive-depictive}

Compare the selected expressive use of \mention{Damn!} with a sound depiction
\mention{Bang!} in a narrative. The bearer in each case is the conventional
use, not every occurrence of the spelling. The question is what contribution
the use makes to interpretation. The strongest simple alternative groups
both as expressive interjections on the basis of their vividness and
standalone packaging. That alternative must still explain the difference
between conveying an attitude and presenting a sound.

A \term{depictive sign} conventionally presents a sound, event, or sensory
scene through a resemblance between aspects of its performance and what
it depicts. Dingemanse's account of ideophones supplies the wider comparison:
conventional lexical items can depict sensory imagery, and their grammatical
organization varies across languages \citep[2--3]{dingemanse2018}.
The term \term{expressive} is itself used for ideophones in some traditions;
that usage doesn't identify them with Potts's expressives.

For the selected \mention{Bang!}, the relevant contribution is a portrayal
of an impact sound. For \mention{Damn!}, it is the conveyed stance toward
a situation. The distinction changes the interpretation of a possible
modification. In a constructed narrative, \mention{bang, bang, bang} can
portray several impacts; repeated \mention{damn} can intensify a stance
without depicting a sequence of impacts. These examples articulate the
contrast to be explained; they aren't results from a new elicitation study.

The literature motivates those different targets. Potts treats repetition
as potentially strengthening expressive contribution \citep[180--182]{potts2007}.
Dingemanse connects depictive performance with iconic mappings between
form and scene \citep[9--11]{dingemanse2018}. An investigation can therefore
ask whether repetition changes portrayed event structure, conveyed stance,
or both. It needn't assume one answer for every form traditionally
called an interjection.

The descriptive comparison has a definite result. Separate expressive
and depictive groupings preserve a difference in contribution that the
undifferentiated alternative omits. A richer alternative could keep one
superordinate class and record the difference as subtypes or dimensions.
For this descriptive task, that is an acceptable equivalent if it retains
the two relations. The result requires the distinction; it doesn't require
a particular taxonomic rank.

Projective category standing remains a further question. Identifying a use
as depictive from its interpretation under repetition and then predicting
that interpretation from the label would recycle the admission evidence.
An independent test could identify the use from unrepeated contexts and
examine modifications elsewhere. Existing accounts motivate that comparison;
the two illustrations don't establish a general response to repetition or
a complete partition of semantic signs.

The groupings also needn't be disjoint. A depiction can convey an attitude,
and an expressive performance can exploit resemblance. What matters is
whether the relevant contribution is conventional in the selected use and
what follows from it. Semantics can investigate both relations, or several
more, without assigning each to a separate department or requiring one
category to represent the semantic perspective as a whole.

\subsection{Contribution, encoding, and competing analyses}

Grouping by contribution differs from grouping by how that contribution
is encoded. \textcite{wharton2016expressives} develops a procedural account
of expressive meaning, relating procedures for interpretation to the
showing of emotion. On such an account, a sign can be expressive in its
contribution and procedural in its encoding. Expressive and procedural
therefore aren't automatically complementary classes.

Three situations need separating. Expressive and depictive groupings can
distinguish contributions within one inquiry. A grouping by contribution
and another by mode of encoding can crosscut one another. But two accounts
making incompatible claims about how the same sign encodes its meaning
compete. Inquiry relativity doesn't adjudicate that competition; evidence
bearing on their divergent consequences must do so.

An inquiry can also support finer groupings. \textcite[113]{ameka1992}
distinguishes emotive and cognitive subtypes within his expressive class:
one concerns feelings or sensations, the other knowledge or recognition.
Those may be useful subtypes for a particular question. Their place in
an analysis depends on the distinctions and generalizations being pursued,
rather than on a requirement that every semantic contrast occupy the same
level of a taxonomy.

\subsection{Interactional resource classes}

Continuers, repair initiators, and closing formats answer different
interactional questions. A continuer such as \mention{mm-hm} can support
another speaker's ongoing turn. Repair \mention{huh?} makes a response to
a problem of hearing or understanding relevant. Closing \mention{bye}
participates in ending an encounter. Accounts of interjections in
conversation make these sequential differences central
\citep[257--259, 263--265]{dingemanse2024}.

For comparison, these can be collected as conventional action formats:
non-elliptical resources that can perform their action without a lexically
specified host. This is a working class, not an assertion that its members
share one next-turn prediction. Such predictions belong initially to
particular resource classes and environments. The contribution of a
superordinate interactional category would need separate support.

The criterion has consequences at the boundary. Response \mention{okay}
and closing or thanking \mention{thanks} are candidates insofar as their
selected uses conventionally perform the relevant action. A formulaic
greeting \mention{how are you?} can qualify despite having several words.
An ordinary imperative \mention{stop!} doesn't qualify merely because
clause structure permits it to stand alone; a claim that it is also a
conventional routine needs evidence beyond that structure.

Host-dependent \mention{please} in \mention{please sit down} belongs to
a different comparison. A standalone appeal \mention{Please!} would have
to be assessed as that selected use. Thus the admission rule concerns
conventional action formats, rather than brevity, wordhood, or every
expression capable of being uttered alone. These constructed comparisons
specify the proposal; they don't supply an annotated inventory.

The grammatical treatment of \mention{mm-hm} and \mention{huh?} likewise
needs more than standalone occurrence. Ameka includes the former among
phatic interjections, and Wilkins treats repair \mention{huh?} as a
conventional utterance \citep[114]{ameka1992}\citep[152]{wilkins1992}.
For the selected uses, I adopt the non-elliptical analysis: neither needs
a recovered nominal head or clause to supply its interpretation.
The grammatical assignments below are defeasible analyses supported by
these descriptions, rather than deductions from punctuation.

Routine formulae expose competition within an inquiry.
\textcite[108--110]{ameka1992} separates \mention{bye}, \mention{sorry},
and \mention{thank you} from interjections, whereas
\textcite[142]{wilkins1992} treats routines as a subtype. Both attend
to meaning and interaction. Assigning each a disciplinary label would
misdescribe their disagreement. The selected closing format can be
compared with continuers without settling whether the wider grouping
is the most useful category for either question.

\subsection{Common bearers without a compulsory matrix}

A \term{sign-type} is an identifiable expression in a conventional use.
Its \term{profile} is the combination of grammatical, semantic, interactional,
and other properties attributed to it. A use can retain its identity
while a proposed property turns out to be absent. Otherwise every
surprising observation would identify a different object and leave the
original proposal immune to correction.

The recurrent contrast between standalone and prenominal \mention{damn}
supports distinguishing their conventional uses. In the complement-taking
case, Table~\ref{tab:contrasts} classifies the lexical head \mention{damn}
in that construction. The whole phrase's external distribution is evidence
about its environment, not a replacement bearer. The table records
analytical judgments, not counts or results of a new annotation study.

\begin{table}[htbp]
\centering
\small
\caption{Selected conventional uses and the relations preserved by named analyses.
The columns organize descriptions; they don't allocate one category to each field.}
\label{tab:contrasts}
\begin{tabular}{@{}>{\raggedright\arraybackslash}p{.24\textwidth}
                  >{\raggedright\arraybackslash}p{.31\textwidth}
                  >{\raggedright\arraybackslash}p{.37\textwidth}@{}}
\toprule
Selected bearer & Grammatical analysis & Meaning or action under comparison \\
\midrule
\mention{gee!}; stance \mention{damn!} & Interjection & Expressive contribution; stance display \\
\addlinespace
\mention{damn} in \mention{damn + NP} & Interjection head or verb: unresolved & Expressive contribution; the head's sequential affordances remain to be specified \\
\addlinespace
Prenominal \mention{damn} & Integrated modifier & Expressive contribution across a grammatical difference \\
\addlinespace
Continuer \mention{mm-hm} & Interjection analysis adopted & Continuation resource; no evaluative stance required by the selected reading \\
\addlinespace
Repair \mention{huh?} & Interjection analysis adopted & Open repair initiator; no evaluative stance required by the selected reading \\
\addlinespace
Closing \mention{bye} & Formulaic standalone expression & Closing resource \\
\addlinespace
Depictive \mention{Bang!} & Interjection or nominal-fragment analysis: unresolved & Sound depiction; stance is a separate question \\
\bottomrule
\end{tabular}
\end{table}

The grammatical and broad interactional assignments can coincide over
selected cases. That coincidence is no defect, and it supplies no
extensional separation between those groupings. The examples do show
expressive commonality across a grammatical difference and different
meaning contributions within a semantic comparison. Their value lies
in the relations they preserve, not in filling every cell of a matrix.

\mention{Bang!} makes the unresolved grammatical question concrete.
English also has nominal \mention{bang}, as in the constructed
\mention{a loud bang} and \mention{two bangs}. A standalone occurrence
therefore needs assessment against a nominal-fragment analysis.
Its depictive interpretation can be recognized while that question remains
open. Section~\ref{sec:phonology} completes the related assessment of
what phonological evidence can decide.

